\chapter*{\textit{Introduction}}\addcontentsline{toc}{chapter}{\textit{Introduction}}
Toute donnée informatique, qu'il s'agisse d'un programme, d'une base de données ou d'un simple document texte, doit être stockée sur un support physique et rendue accessible de manière structurée. Cette organisation logique du stockage est assurée par le \textbf{système de fichiers} (\textit{File System}), composant fondamental de tout système d'exploitation moderne.

Un système de fichiers définit la manière dont les données sont organisées, nommées, indexées et retrouvées sur les mémoires secondaires~: disques durs, clés USB, CD-ROM ou disques SSD. Il constitue l'interface entre le matériel physique brut et l'utilisateur, lui offrant une vue abstraite et hiérarchique de ses données à travers un chemin d'accès.

Le présent exposé se propose d'étudier les systèmes de fichiers de manière progressive et structurée. Nous aborderons dans un premier temps les \textbf{généralités} sur les systèmes de fichiers (définition, notion de fichier, représentation pour l'utilisateur), puis leur \textbf{organisation hiérarchique} sous Windows et Unix/Linux. Nous examinerons ensuite le mécanisme essentiel de \textbf{montage et démontage}, avant de dresser une classification des \textbf{principaux types de systèmes de fichiers} existants, des systèmes locaux aux systèmes journalisés, en passant par les systèmes réseau et optiques.